\subsection{Exercise 6}
Tích chập tuyến tính: $y(n) = \sum_{k=-\infty}^{\infty} h(k)\,x(n-k)$.
\begin{enumerate}
    \item Rút gọn $y(n)$ để tổng hội tụ trong hai trường hợp:
    \begin{itemize}
        \item $h(n)$ và $x(n)$ cùng nhân quả.
        \item Hệ FIR: $h(n) = 0$ với $n < 0 \lor n > M$.
    \end{itemize}
    \item (Trường hợp 2) Tính $y(n)$ với $x(n) = \{\underset{\uparrow}{1}, 2, -3, 2, 1\}$, $h(n) = \{\underset{\uparrow}{1}, 0, -1\}$ bằng phương pháp gập--dịch. Hiện thực bằng Scilab.
    \item (Trường hợp 2) Tính lại $y(n)$ bằng phương pháp ma trận. Hiện thực bằng Scilab.
    \item Vẽ $x(n)$, $y(n)$ và năng lượng tín hiệu trên cùng một figure.
\end{enumerate}

% TODO: lời giải

% \begin{lstlisting}
% // Scilab code
% \end{lstlisting}

% \begin{figure}[H]
%     \centering
%     \includegraphics[width=0.75\textwidth]{img/chap5/ex6.png}
%     \caption{...}
% \end{figure}
